\documentclass[12pt]{article}
\usepackage[T1]{fontenc}
\usepackage[utf8]{inputenc}
\usepackage{lmodern}
\usepackage{textcomp}
\usepackage{enumitem}
\usepackage{geometry}
\geometry{margin=1in}
\begin{document}
\begin{center}
Answer Key - Astronomy - Graduate Level Assessment 13 \
\small (Answers are ordered exactly as in the question paper.)
\end{center}

\section*{Multiple Choice}
\begin{enumerate}[label=\arabic*.]
\item \textbf{Answer:} D
\\ \textit{Options:} A) Iron; B) Hydrogen; C) Methane; D) Helium
\\ \textit{Note:} Helium is the second most abundant element in the solar system, comprising about 24\% of its elemental composition, primarily produced through nuclear fusion in stars, following hydrogen, which accounts for approximately 74\%.
\item \textbf{Answer:} B
\\ \textit{Options:} A) Being hit in the head by a bullet.; B) Being hit by a small meteorite.; C) Starvation during global winter caused by a major impact.; D) Driving while intoxicated without wearing seatbelts.
\\ \textit{Note:} Being hit by a small meteorite is the least likely cause of death because such events are extremely rare, with the likelihood of an individual being struck by a meteorite being astronomically low compared to other causes of death.
\item \textbf{Answer:} D
\\ \textit{Options:} A) From the layering of materials within the Earth.; B) From fossils of ancient life.; C) From the cratering history of Earth's surface.; D) From radioactive dating of rocks and meteorites.
\\ \textit{Note:} The age of the Earth is determined through radioactive dating, which measures the decay of radioactive isotopes in rocks and meteorites, providing a reliable timeline of Earth's formation and its geological history.
\item \textbf{Answer:} D
\\ \textit{Options:} A) the softer rocky material of the mantle.; B) the lava that comes out of volcanoes.; C) material between the crust and the mantle.; D) the rigid rocky material of the crust and uppermost portion of the mantle.
\\ \textit{Note:} The lithosphere is defined as the rigid outer layer of a planet, encompassing both the solid rock of the crust and the uppermost part of the mantle, which collectively provide structural support and influence geological processes.
\item \textbf{Answer:} D
\\ \textit{Options:} A) patterns of motion; B) the presence of asteroids and comets; C) the measured ages of meteorites; D) the presence of life on Earth
\\ \textit{Note:} A formation theory of the solar system primarily focuses on the processes that led to the formation and arrangement of celestial bodies, while the presence of life on Earth is a separate issue that involves biological, geological, and environmental factors beyond the scope of solar system formation.
\end{enumerate}

\section*{True / False}
\begin{enumerate}[label=\arabic*.]
\setcounter{enumi}{5}
\item \textbf{Answer:} True
\\ \textit{Options:} A) True; B) False
\\ \textit{Note:} The statement is true because one astronomical unit (AU) is standardized as the average distance between the Earth and the Sun, which is approximately 150 million kilometers.
\item \textbf{Answer:} False
\\ \textit{Options:} A) True; B) False
\\ \textit{Note:} Venus experiences no seasons not because of its slow rotation, but because its axis tilt is only about 3 degrees, resulting in minimal variation in solar radiation across its surface throughout its orbit.
\item \textbf{Answer:} True
\\ \textit{Options:} A) True; B) False
\\ \textit{Note:} The Pleiades, also known as the Seven Sisters, is an open star cluster that contains several bright stars, with seven being the most prominent and easily visible to the naked eye, thereby validating the statement as true.
\item \textbf{Answer:} False
\\ \textit{Options:} A) True; B) False
\\ \textit{Note:} Ganymede is actually the farthest of Jupiter's four largest moons, with Io being the closest, followed by Europa, and then Callisto.
\item \textbf{Answer:} False
\\ \textit{Options:} A) True; B) False
\\ \textit{Note:} The statement is false because the rock labeled as 'bigfoot' is actually much farther away than 10 meters, with accurate measurements indicating a distance of several hundred meters from the rover.
\end{enumerate}

\section*{Long Form Response}
\begin{enumerate}[label=\arabic*.]
\setcounter{enumi}{10}
\item \textbf{Answer:} Saturn's comparable size to Jupiter, despite its lower mass, is primarily attributed to differences in density and the effects of gravitational compression. Jupiter's greater mass exerts a stronger gravitational force, leading to increased compression of its gaseous envelope, which results in a higher average density. In contrast, Saturn's lower mass allows for a less intense gravitational pull, resulting in a larger volume but a lower density. This difference in structure highlights the relationship between mass and gravitational forces in shaping planetary characteristics, demonstrating that size does not always correlate directly with mass. Consequently, while both planets share similar dimensions, their distinct compositions and the effects of gravitational compression reveal fundamental differences in their internal structures and physical properties.
\\ \textit{Note:} correct because it accurately describes how Jupiter's greater mass leads to higher gravitational compression and density, resulting in a smaller volume compared to Saturn, which, due to its lower mass, has a lower density and a comparatively larger size.
\item \textbf{Answer:} Mass and weight are closely related yet distinct concepts in physics; mass is a measure of the amount of matter in an object, while weight is the gravitational force acting on that mass. The weight of an object can be calculated using the formula \( W = mg \), where \( W \) is weight, \( m \) is mass, and \( g \) is the acceleration due to gravity. On Earth, the gravitational pull is approximately \( 9.81 \, \text{m/s}^2 \), which means a 5 kg sandbag has a weight of \( 5 \, \text{kg} \times 9.81 \, \text{m/s}^2 = 49.05 \, \text{N} \). However, on a planet with half the gravitational pull of Earth, where \( g \) would be \( 4.905 \, \text{m/s}^2 \), the weight of the same 5 kg sandbag would be \( 5 \, \text{kg} \times 4.905 \, \text{m/s}^2 = 24.525 \, \text{N} \). Thus, while the mass remains constant at 5 kg, the weight is significantly reduced to approximately 25 N due to the lower gravitational force.
\\ \textit{Note:} The answer correctly identifies that mass is a constant property of matter and that weight varies with gravitational force, demonstrating this relationship by calculating the weight of a 5 kg sandbag on a planet with half of Earth's gravitational pull, thus highlighting the fundamental dependence of weight on gravity.
\end{enumerate}

\vfill
\noindent\textit{End of Answer Key}
\end{document}